\section{The Challenge}

Partial Pre/Post-Conditions + C Tests to Formal API Semantics

\subsection{Example: Integer Div/Mod}

Laws of mod/div for integers in programming languages (e.g. Haskell):

\begin{eqnarray*}
   n \div (-d)    &=& -(n \div d)
\\ (-n) \div d    &=& -(n \div d)
\\ (-n) \div (-d) &=& n \div d
\end{eqnarray*}

\begin{eqnarray*}
   n \mod (-d)    &=& -(n \mod d)
\\ (-n) \mod d    &=& n \mod d
\\ (-n) \mod (-d) &=& -(n \mod d) 
\end{eqnarray*}

\def\abs#1{|{\!#1\!}|}
\def\sgn#1{\mathsf{sgn}~{#1}}

To reason about the above we also need functions for 
sign ($\sgn\_$)
and absolute values ($\abs{\_}$).

\begin{eqnarray*}
   \sgn n  &=& 1, \qquad n > 0
\\         &=& -1, \qquad n < 0
\\         &=& 0, \qquad \text{otherwise}
\end{eqnarray*}

\begin{eqnarray*}
   \abs n  &=& -n, \qquad n < 0
\\         &=& n, \qquad \text{otherwise}
\end{eqnarray*}

We assume that we know about integer arithmetic and comparisons.

We have preconditions regarding the ``zeroness`` of the divisor,
and the sign of both inputs. The sign-function can establish these.
We can express postconditions in terms of 
the absolute values and signs of the inputs.

\newpage
\subsection{Preconditions}

Those regarding $n$:
\begin{eqnarray*}
   N0 && \sgn n = 0
\\ NP && \sgn n = 1
\\ NN && \sgn n = -1
\end{eqnarray*}

Those regarding $d$:
\begin{eqnarray*}
   D0 && \sgn d = 0
\\ DP && \sgn d = 1
\\ DN && \sgn d = -1
\end{eqnarray*}

\subsection{Transitions}

We simplfiy by assuming functions $\div'$ and $\mod'$
that work correctly on positive numbers 

Starting with $\div$:
\begin{eqnarray*}
   n \div d &=& r
\\ \_,D0 && r = \bot \qquad  UNDEFDIV
\\ NP,DP && r = \abs n \div' \abs d \qquad  ISDIV
\\ NN,DP && r = -(\abs n \div' \abs d) \qquad  NEGDIV
\\ NP,DN && r = -(\abs n \div' \abs d) \qquad  NEGDIV
\\ NN,DN && r = \abs n \div' \abs d\qquad  ISDIV
\end{eqnarray*}

Now, $\mod$:
\begin{eqnarray*}
   n \mod d &=& r
\\ \_,D0 && r = \bot \qquad  UNDEFMOD
\\ NP,DP && r = \abs n \mod' \abs d \qquad ISMOD
\\ NN,DP && r = \abs n \mod' \abs d \qquad ISMOD
\\ NP,DN && r = -(\abs n \mod' \abs d) \qquad NEGMOD
\\ NN,DN && r = -(\abs n \mod' \abs d) \qquad NEGMOD
\end{eqnarray*}

\subsection{Postconditions}

UNDEFDIV to NEGMOD